\scanpage{5}
% Lines 1–2, continuation of the final line of scan 3.
\margin{$\Vert$}sont duels l'un de l'autre, et que $H^p(A,\Omega_A^q)$ est dual de $H^q(B,\Omega_A^p)$.

% Lines 3–9. There is an opening square bracket in the left margin.
[ Ce dernier fait, grâce à $H^*(A,\Omega_A^*)=\bigwedge\bigl(H^*(A,\Omega_A^{*})\bigr)^{(1)}$ se ramène au cas de la \unsure{dim 1|dim 1.}, i.e. par symétrie à la dualité de $H^{\deleted{\unsure{0}}1}(A,\deleted{\unsure{\Omega}}\underline{O}_A)$ et $H^0(B,\deleted{\unsure{\underline{O}}}\Omega_B^1)$, qui est immédiate. \deleted{De plus, \unsure{on},} pour avoir une dualité entre $H^*(A)$ et $H^*(B)$, il \unsure{suffit|\deleted{suffit}|suffirait|\deleted{suffirait}} grâce à la structure connue de la définir entre $H^1(A)$ et $H^1(B)$, ce qui ne saurait offrir de difficultés....

% Lines 10–15
Il y a lieu enfin de comparer ces extensions et dualités avec ce que donne la \uline{dualité de Cartier}. \unsure{Envisageons|Considérons} (\unsure{la car.|le cas} étant $p>0$) $ {}_pA$ qui est dual au sens Cartier de $ {}_pB$. On a aussi $ {}_F A=\operatorname{Gr}(t_A)\subset {}_p A$\unsure{.|-}\,; de degrés respectifs $p^{\deleted{\unsure{2|2d}}d}$ et $p^{2d}$, où $d=\dim A$.
% Lines 16–17
\[
 \left.\begin{aligned}
 {}_F A&\subset {}_p A\\
 {}_F B&\subset {}_p B
 \end{aligned}\right]\ \text{duels l'un de l'autre}
\]
% Lines 18–19
Je dis que $ {}_F A$ \uline{et} $ {}_F B$ \uline{sont orthogonaux l'un de l'autre}. \deleted{Noter que $ {}_F B$} On le
% archive sheet number: 4
